% A3: Warmup Prior Transfer — Acceleration, Risk, and Limitations
% Analyzes when semantic transfer helps vs. hurts, and how Corralling mitigates risk.
% Cross-references: D1 (warmup vs tabula rasa, empirical validation), A1 (composite bound)

\subsection{Warmup Prior Transfer: Acceleration, Risk, and Limitations}
\label{appendix:warmup_transfer}

The banditGPT warmup expert initializes its LinUCB matrices using priors
derived from 80k LMSYS Chatbot Arena battles.  This subsection analyzes
(i) when this transfer accelerates convergence, (ii) the risk when priors are
misspecified, and (iii) how the Corralling architecture bounds that risk.  We
also discuss known limitations of the current formulation.

\subsubsection{Prior Initialization and Effective Sample Size}

Under the default Hybrid LinUCB policy (Algorithm~\ref{alg:linucb}),
initialization involves both arm-specific and family-shared state.

\paragraph{Arm-specific initialization.}
For a new or existing model~$a$, the warmup expert initializes:
\begin{align}
A_a^{(0)} &= \neff \cdot \lambda \cdot I_d \label{eq:warmup_A} \\
b_a^{(0)} &= \neff \cdot \lambda \cdot \theta_{\text{neighbor}} \label{eq:warmup_b}
\end{align}
where $\theta_{\text{neighbor}}$ is the learned reward parameter of the most
similar model in embedding space (identified via cosine similarity of model
name embeddings), $\neff$ is the effective prior sample size, and $\lambda$ is
the ridge regularization parameter.

This yields the initial estimate $\hat{\theta}_a = (A_a^{(0)})^{-1} b_a^{(0)} =
\theta_{\text{neighbor}}$ with posterior variance proportional to $1 / \neff$.
The Bayesian interpretation is straightforward: the system behaves as if it has
already observed $\neff$ samples from the new model, all consistent with the
neighbor's learned behavior.

\paragraph{Family-shared initialization.}
For each family $g \in \mathcal{G}$, the family-shared matrices start from an
uninformative prior:
\begin{align}
A_g^{(0)} &= \lambda \cdot I_d \label{eq:warmup_A0} \\
b_g^{(0)} &= \mathbf{0} \label{eq:warmup_b0}
\end{align}
The family-shared component $\beta_g$ is learned entirely online: as any arm
$a$ with $F(a) = g$ receives observations, both $A_g^{(0)}$ and $b_g^{(0)}$
are updated with the \emph{full} reward signal, while the arm-specific
$A_a, b_a$ learn only the residual $r_t - x_t^\top \hat{\beta}_g$.  When a
new model joins an existing family via \texttt{register\_model()}, the
already-trained $\beta_g$ provides immediate calibration---the model
onboarding experiment (Section~\ref{sec:results}) shows 46\% lower initial
prediction error compared to isolated-family assignment.

\subsubsection{Acceleration Under Accurate Priors}

\begin{definition}[Prior Accuracy]
\label{def:prior_accuracy}
A warmup prior $\theta_{\text{neighbor}}$ is \emph{$\epsilon$-accurate} for model $a$ if
$\|\theta_{\text{neighbor}} - \theta_a^*\| \leq \epsilon$, where $\theta_a^*$
is the true reward parameter.
\end{definition}

When the prior is $\epsilon$-accurate, the warmup expert effectively begins
with a confidence set of radius ${\sim}\epsilon$ around the true parameter.  A
tabula rasa expert starting from $\hat{\theta} = 0$ must explore until its
confidence set shrinks to comparable radius, which requires $\Omega(d / \epsilon^2)$
observations under standard LinUCB analysis~\cite{abbasi2011improved}.

\begin{theorem}[Informal: Warmup Acceleration]
\label{thm:warmup_acceleration}
If the warmup prior is $\epsilon$-accurate for all arms, the warmup expert's
regret satisfies:
\begin{equation}
R_{\textnormal{warmup}}(T)
\;\lessapprox\;
R_{\textnormal{tabula rasa}}(T) - R_{\textnormal{tabula rasa}}(\neff)
\label{eq:warmup_savings}
\end{equation}
That is, the warmup expert ``skips'' approximately $\neff$ rounds of
exploration, avoiding the regret incurred during the cold-start phase.
\end{theorem}

\paragraph{Empirical confirmation.}
The Corralling ablation (Section~\ref{sec:appendix_corralling_ablation}) uses
a 3-model portfolio where GPT-4o is initialized via semantic transfer from
GPT-4-Turbo.  Despite having no direct training data, GPT-4o is discovered as
the dominant choice (64.6\% usage at $\lambda=0$), demonstrating that semantic
transfer provides effective cold-start acceleration on real LMSYS Arena data.

\subsubsection{Risk Under Misspecified Priors}

When $\theta_{\text{neighbor}}$ is far from $\theta_a^*$ (domain mismatch,
model capability differences, or outdated priors), the warmup expert incurs
additional regret from \emph{negative transfer}.

\begin{definition}[Prior Misspecification]
\label{def:misspecification}
A warmup prior is \emph{$\delta$-misspecified} for model $a$ if
$\|\theta_{\text{neighbor}} - \theta_a^*\| \geq \delta$ for some $\delta > \epsilon$.
\end{definition}

Under misspecification, the warmup expert's initial confidence set is centered
at a point that is $\delta$-far from the truth.  The expert must accumulate
sufficient data to ``overwrite'' the prior, requiring $\Omega(\neff)$
observations before the posterior mean moves substantially toward $\theta_a^*$.
During this period, the expert makes systematically biased selections.

\begin{remark}[The $\neff$ dilemma]
\label{remark:neff_dilemma}
Higher $\neff$ accelerates convergence when the prior is accurate but
\emph{slows} recovery when it is misspecified: the system must see more data to
override a more confident (but wrong) prior.  This creates a tension that
cannot be resolved by tuning $\neff$ alone, because the quality of the
semantic neighbor is not known in advance.
\end{remark}

\subsubsection{How Corralling Bounds Negative Transfer}

The composite regret decomposition (Theorem~\ref{thm:composite}) provides the
key insight: the system's regret tracks the \emph{better} expert, not the
worse one.

\begin{corollary}[Bounded Negative Transfer]
\label{cor:neg_transfer}
Even when the warmup prior is arbitrarily misspecified, the total system regret
satisfies:
\begin{equation}
R_{\textnormal{total}}(T)
\;\leq\;
R_{\textnormal{tabula rasa}}(T)
\;+\;
O\!\bigl(\sqrt{T \ln 2}\bigr)
\label{eq:neg_transfer_bound}
\end{equation}
That is, the worst case is clean-slate LinUCB performance plus the meta-learner
overhead---not the arbitrarily bad performance of the misspecified warmup expert.
\end{corollary}

\begin{proof}
Immediate from Theorem~\ref{thm:composite}: $R_{\text{total}}(T) \leq
O(\sqrt{T \ln 2}) + \min\{R_1(T), R_2(T)\} \leq O(\sqrt{T \ln 2}) + R_2(T)$.
\end{proof}

\paragraph{The cost of safety.}
The $O(\sqrt{T \ln 2})$ term is the price paid for not committing to either
expert up front.  This is the central design claim of the banditGPT
architecture:

\begin{quote}
\emph{Corralling converts the binary risk of warmup priors---helpful or
harmful---into a bounded $O(\sqrt{T})$ overhead, while preserving the full
acceleration when priors are accurate.}
\end{quote}

\subsubsection{Correct vs.\ Naive Prior Injection}
\label{appendix:naive_vs_correct}

Our Bayesian formulation
(Equations~\eqref{eq:warmup_A}--\eqref{eq:warmup_b}) scales \emph{both} the
arm-specific precision matrix and moment vector by $\neff$, preserving the
mean prediction $\hat{\theta}_a = \theta_{\text{neighbor}}$ regardless of
$\neff$.  The family-shared matrices
(Equations~\eqref{eq:warmup_A0}--\eqref{eq:warmup_b0}) are deliberately
\emph{not} warm-started: $\beta_g$ is learned purely online, so the warmup
prior affects only the arm-specific residual channel.
A naive implementation that scales only the moment vector:
\begin{equation}
A_{\text{naive}} = \lambda I, \quad b_{\text{naive}} = \neff \cdot \lambda \cdot \theta_{\text{neighbor}}
\end{equation}
produces a mean prediction $\hat{\theta}_{\text{naive}} = \neff \cdot
\theta_{\text{neighbor}}$, which artificially inflates predicted rewards by a
factor of $\neff$.  This ``optimistic bias'' forces selection of the new model
regardless of context.  The correct formulation ensures that $\neff$ controls
only \emph{confidence} (variance), not the \emph{direction} of the prior.

\paragraph{Practical $\neff$ recommendation.}
Since $\neff$ controls variance but not mean, the choice is forgiving when the
neighbor is a good match.  When neighbor quality is uncertain, moderate values
($\neff \in [2, 10]$) hedge against both under-confidence (slow convergence)
and over-confidence (slow recovery from mismatch).  The Corralling safety net
(Corollary~\ref{cor:neg_transfer}) provides a second layer of protection: even
with a high $\neff$ and a poor neighbor, total regret is bounded by
tabula-rasa performance plus $O(\sqrt{T})$ meta-overhead.

\subsubsection{Known Limitations}
\label{appendix:warmup_limitations}

We identify several limitations of the current formulation that represent
directions for future work.

\paragraph{1. Isotropic regularization in PCA space.}
The prior matrices use isotropic initialization $A_a^{(0)} = \neff \lambda I$,
which regularizes all PCA dimensions equally.  After PCA, principal components
have decreasing empirical variance, so equal regularization
over-shrinks low-variance components relative to their scale.  An anisotropic
regularizer $A_a^{(0)} = \neff \lambda \cdot \text{diag}(\sigma_1^{-2},
\ldots, \sigma_d^{-2})$ that scales with per-component variance would be more
principled but adds implementation complexity and requires access to the PCA
eigenspectrum at initialization time.

\paragraph{2. $\neff$ conflates prior strength with prior accuracy.}
The current system uses a single scalar $\neff$ that controls both the
confidence (variance) and influence (number of pseudo-observations) of the
prior.  Ideally, these would be decoupled: a system could express ``I am
confident this prior is directionally correct but uncertain about its
magnitude'' by using a high $\neff$ for the mean direction and a low
$\neff$ for the scale.  The current implementation does not support this
decomposition.

\paragraph{3. Stationary expert-level updates.}
Each LinUCB expert accumulates observations monotonically (the precision
matrix $A$ only grows), which is appropriate under stationary reward
distributions but cannot track temporal drift in model capabilities.
Non-stationarity is handled only at the meta level (Corralling shifts weight
between experts) and through the production library's optional forgetting
factor ($\gamma_{\text{decay}} < 1$, disabled in experiments).  Extending
expert-level updates with explicit change-point detection or sliding-window
regression is a direction for future work.

\paragraph{4. Linear realizability.}
The regret guarantees in Section~\ref{appendix:regret_decomposition} assume
linear reward structure (Equation~\ref{eq:linear_reward}).  Prompt--model
reward surfaces may exhibit non-linear interactions (e.g., a model that excels
at short creative prompts but struggles with long creative prompts).  Hybrid
LinUCB provides per-arm linear fits (via the combined
$\beta_g + \theta_a$ prediction), which offer piecewise flexibility across
models and within-family transfer, but cannot capture within-arm
non-linearity.  Kernel or neural extensions~\cite{zhou2020neural} could
address this at the cost of computational overhead.

\paragraph{5. Semantic neighbor selection.}
The quality of warmup transfer depends critically on the semantic neighbor
identified via model name embedding similarity.  This is a heuristic: two
models with similar names (e.g., ``GPT-4-Turbo'' and ``GPT-4o'') may have
substantially different reward profiles.  The Corralling safety net
(Corollary~\ref{cor:neg_transfer}) bounds the damage from poor neighbors, but
a more sophisticated neighbor selection mechanism (e.g., using task-specific
evaluation data or capability metadata) could improve transfer quality
directly.
